% Part: incompleteness
% Chapter: incompleteness-provability
% Section: rosser-thm

\documentclass[../../../include/open-logic-section]{subfiles}

\begin{document}

\olfileid[fr]{inc}{inp}{ros}

\olsection{Le théorème de Rosser}

Peut-on renforcer la preuve de Gödel en remplaçant
la « $\omega$-cohérence » par la seule « cohérence »?
Oui, grâce à une astuce de Rosser : employer un
prédicat de démontrabilité « modifié »
$\ORProv_T(y)$ à la place de
$\OProv[\Th{T}](y)$.

\begin{thm}
\ollabel{thm:rosser}
Toute théorie~$\Th{T}$ cohérente,
!!{axiomatizable} et étendant~$\Th{Q}$ est
incomplète.
\end{thm}

\begin{proof}
Rappelons que $\OProv[\Th{T}](y)$ est défini par
$\lexists[x][\OPrf[\Th{T}](x,y)]$ : la
!!{formula}~$\OPrf[\Th{T}](x,y)$ représente la
relation décidable selon laquelle $x$ est le numéro
de Gödel d'une !!{derivation}, dans~$\Th{T}$,
de l'!!{sentence} de numéro~$y$.
La relation qui dit que $x$ code une
\emph{réfutation} de cet !!{sentence} est elle
aussi décidable. Soit $\fn{not}(x)$ la fonction
récursive primitive qui, si $x$ code une
!!{formula}~$!A$, renvoie le code de~$\lnot !A$; sinon, elle renvoie $0$.
Alors $\Refut[\Th{T}](x,y)$ vaut si et seulement
si $\Prf[\Th{T}](x,\fn{not}(y))$ vaut.
Notons $\ORefut[\Th{T}](x,y)$ une !!{formula}
qui représente cette relation. Si
$\Th{T}\Proves\lnot !A$ et si $\delta$ est une
!!{derivation} correspondante, alors
$\Th{Q}\Proves\ORefut[\Th{T}](\gn{\delta},\gn{!A})$.
Définissons maintenant $\ORProv[\Th{T}](y)$ par
\[
\lexists[x][(\OPrf[\Th{T}](x,y)\land
\lforall[z][(z<x\lif\lnot\ORefut[\Th{T}](z,y))])].
\]
Cette !!{formula} affirme approximativement qu'il
existe une preuve de l'!!{sentence} de code~$y$
dans~$\Th{T}$, mais aucune réfutation de code
plus petit. Si $\Th{T}$ est cohérente,
$\ORProv[\Th{T}](y)$ et $\OProv[\Th{T}](y)$
sont vraies des mêmes nombres. Du point de vue de
ce que $\Th{T}$ peut \emph{démontrer}, leurs
propriétés diffèrent pourtant : vérité et
démontrabilité ne se confondent pas. Si $\Th{T}$
est \emph{incohérente}, la coïncidence des deux
prédicats n'est plus garantie. On lit parfois
$\ORProv[\Th{T}](y)$ « $y$ est démontrable au
sens de Rosser ». Cette expression pourrait laisser
croire qu'il s'agit d'une sorte particulière de
démontrabilité; or, dans une théorie incohérente,
certains !!{sentence}s peuvent être démontrables sans
l'être au sens de Rosser. Pour éviter cette
confusion, on pourrait lui donner un nom
arbitraire, tel que « shmovable ».

Par le lemme du point fixe, il existe un
!!{sentence}~$!R_\Th{T}$ tel que
\begin{equation}
\Th{Q}\Proves !R_\Th{T}\liff
\lnot\ORProv[\Th{T}](\gn{!R_\Th{T}}).
\ollabel{RT}
\end{equation}
Contrairement à la preuve du premier théorème
d'incomplétude (\olref[1in]{thm:first-incompleteness}),
montrons qu'une théorie cohérente~$\Th{T}$ ne démontre ni
$!R_\Th{T}$ ni $\lnot !R_\Th{T}$, sans supposer
la $\omega$-cohérence.

Montrons d'abord que $\Th{T}\Proves/ !R_{\Th{T}}$.
Supposons le contraire et soit $n$ le numéro de
Gödel d'une !!{derivation} de~$!R_{\Th{T}}$
à partir de~$T$. Puisque $\OPrf[\Th{T}]$
représente $\Prf[\Th{T}]$ dans~$\Th{Q}$,
\[
\Th{Q}\Proves
\OPrf[\Th{T}](\num{n},\gn{!R_{\Th{T}}}).
\]
Comme $\Th{T}$ est cohérente, aucun entier
$k<n$ ne code une !!{derivation} de
$\lnot !R_{\Th{T}}$. Pour chacun de ces~$k$,
$\Th{Q}\Proves\lnot\ORefut[\Th{T}](\num{k},
\gn{!R_{\Th{T}}})$. La borne finie sur les
entiers naturels
(\olref[req][min]{lem:less-nsucc}) permet d'en
déduire
$\Th{Q}\Proves\lforall[z][(z<\num{n}\lif
\lnot\ORefut[\Th{T}](z,\gn{!R_{\Th{T}}}))]$.
Ainsi,
\[
\Th{Q}\Proves
\lexists[x][(\OPrf[\Th{T}](x,\gn{!R_{\Th{T}}})
\land\lforall[z][(z<x\lif
\lnot\ORefut[\Th{T}](z,\gn{!R_{\Th{T}}}))])],
\]
c'est-à-dire $\ORProv[\Th{T}](\gn{!R_{\Th{T}}})$.
L'équivalence ci-dessus (\olref{RT}) donne alors
$\Th{Q}\Proves\lnot !R_{\Th{T}}$; puisque
$\Th{T}$ étend~$\Th{Q}$, elle démontre aussi
cette négation. Or nous avons supposé qu'elle
démontre $!R_{\Th{T}}$. Sa cohérence serait
contredite.

Montrons ensuite que
$\Th{T}\Proves/\lnot !R_{\Th{T}}$. Supposons
qu'il existe une !!{derivation} de cette négation
de numéro de Gödel~$n$. Alors
$\Refut[\Th{T}](n,\Gn{!R_{\Th{T}}})$ vaut et,
par représentation,
$\Th{Q}\Proves\ORefut[\Th{T}](\num{n},
\gn{!R_{\Th{T}}})$. Nous allons montrer que
$\Th{T}$ démontrerait aussi $!R_{\Th{T}}$ et
serait incohérente. Comme
\begin{align*}
\Th{Q} & \Proves !R_{\Th{T}}\liff
\lnot\ORProv[\Th{T}](\gn{!R_{\Th{T}}}),
\intertext{et que $\Th{T}$ étend $\Th{Q}$, il suffit d'établir}
\Th{Q} & \Proves
\lnot\ORProv[\Th{T}](\gn{!R_{\Th{T}}}).
\end{align*}
L'!!{sentence} $\lnot\ORProv[\Th{T}]
(\gn{!R_{\Th{T}}})$, c'est-à-dire
\begin{align*}
\lnot & \lexists[x][(\OPrf[\Th{T}](x,
\gn{!R_{\Th{T}}})\land
\lforall[z][(z<x\lif
\lnot\ORefut[\Th{T}](z,\gn{!R_{\Th{T}}}))])],
\intertext{équivaut logiquement à}
& \lforall[x][(\OPrf[\Th{T}](x,
\gn{!R_{\Th{T}}})\lif
\lexists[z][(z<x\land
\ORefut[\Th{T}](z,\gn{!R_{\Th{T}}}))])].
\end{align*}
Raisonnons informellement à l'intérieur de
$\Th{Q}$, en utilisant ce qu'elle démontre.
Soit $x$ quelconque et supposons
$\OPrf[\Th{T}](x,\gn{!R_{\Th{T}}})$.
Nous savons déjà que
$\Th{T}\Proves/ !R_{\Th{T}}$. Pour chaque entier
fixé~$k$, la relation
$\Prf[\Th{T}](k,\Gn{!R_{\Th{T}}})$ est donc
fausse et, par représentation négative,
$\Th{Q}\Proves\lnot\OPrf[\Th{T}](\num{k},
\gn{!R_{\Th{T}}})$. Sous notre hypothèse sur~$x$,
on en déduit $\eq/[x][\num{k}]$.
Nous n'utiliserons que les cas finis $k=0,\dots,n$.
En particulier, $x\ne\num{n}$; les cas
$k<n$ donnent aussi
$\lnot(\eq[x][\num{0}]\lor\eq[x][\num{1}]
\lor\dots\lor\eq[x][\num{n-1}])$.
La propriété des bornes finies
(\olref[req][min]{lem:less-nsucc}) en tire
$\lnot(x<\num{n})$. Avec la trichotomie
(\olref[req][min]{lem:trichotomy}), il s'ensuit
$\num{n}<x$. Puisque
$\Th{Q}\Proves\ORefut[\Th{T}](\num{n},
\gn{!R_{\Th{T}}})$, nous obtenons
$\num{n}<x\land\ORefut[\Th{T}](\num{n},
\gn{!R_{\Th{T}}})$, puis
$\lexists[z][(z<x\land
\ORefut[\Th{T}](z,\gn{!R_{\Th{T}}}))]$.
Comme $x$ était quelconque, nous avons donc
\[
\lforall[x][(\OPrf[\Th{T}](x,
\gn{!R_{\Th{T}}})\lif
\lexists[z][(z<x\land
\ORefut[\Th{T}](z,\gn{!R_{\Th{T}}}))])].
\]
C'est l'énoncé voulu.
\end{proof}

\begin{prob}
Deux ensembles $A$ et $B$ d'entiers naturels sont
\emph{effectivement inséparables} s'il n'existe
aucun ensemble !!{decidable}~$X$ tel que
$A\subseteq X$ et $B\subseteq\Complement{X}$,
où $\Complement{X}=\Nat\setminus X$.
Soit $\Th{T}$ une extension cohérente et
!!{axiomatizable} de~$\Th{Q}$. Supposons que $A$
soit l'ensemble des numéros de Gödel des
!!{sentence}s démontrables dans~$\Th{T}$, et $B$
celui des numéros de Gödel des !!{sentence}s
réfutables dans~$\Th{T}$. Montrez que $A$ et~$B$
sont effectivement inséparables.
\end{prob}

\end{document}
